\documentclass[12pt,a4paper]{article}
\usepackage[utf8]{inputenc}
\usepackage[T1]{fontenc}
\usepackage{amsmath,amssymb,amsthm}

\theoremstyle{plain}
\newtheorem{theorem}{Theorem}
\newtheorem{proposition}[theorem]{Proposition}
\newtheorem{corollary}[theorem]{Corollary}

\title{A Phase Symmetry Conjecture for the Riemann Zeta Function}
\author{Andr\'e Gualberto}

\begin{document}
\maketitle

\begin{abstract}
We examine the condition $K(s)K(1-s)\in\mathbb{R}$ for
$K(s)=2^s-1$, which is algebraically equivalent to
$\Re(s)=1/2$ or $\Im(s)=n\pi/\ln2$ for integer $n$.
If a non-trivial zero of the Riemann zeta function were to
lie off the critical line, it could only do so on one of
these exceptional lines.  We present a numerical scan of
200 such lines that finds no candidate.  The exclusion of
these lines remains an open problem.
\end{abstract}

\section{Introduction}

The fundamental equation
\[
\boxed{\tfrac12 \times 2 = 1}
\]
is the algebraic seed of the symmetry explored here.  The number
$2$ is the base of the parity alternation that separates odd primes
from even numbers; $1/2$ is the critical line of the Riemann
hypothesis; and $1$ is the primordial unit.

The valence function $K(s)=2^s-1$ captures the role of the
base~$2$ in the functional equation of $\zeta(s)$.
The condition $K(s)K(1-s)\in\mathbb{R}$ forces
$\Re(s)=1/2$ unless $\Im(s)=n\pi/\ln2$.

\section{Algebraic Symmetry}

Define $K(s)=2^s-1$.

\begin{proposition}
\label{prop:symmetry}
$K(s)K(1-s)\in\mathbb{R}$ iff $\Re(s)=\tfrac12$ or
$\Im(s)=n\pi/\ln2$ for some $n\in\mathbb{Z}$.
\end{proposition}
\begin{proof}
Write $s=\sigma+it$.  Then
\[
K(s)K(1-s) = (2^s-1)(2^{1-s}-1) = 3 - 2^s - 2^{1-s}.
\]
Since $2^{1-s}=2\cdot2^{-s}$,
\[
\Im\bigl(K(s)K(1-s)\bigr) = -\Im(2^s + 2^{1-s}).
\]
With $2^s = e^{\sigma\ln2}\,e^{it\ln2}$,
\begin{align*}
\Im(2^s + 2^{1-s})
&= e^{\sigma\ln2}\sin(t\ln2)
   + e^{(1-\sigma)\ln2}\sin(-t\ln2)\\
&= \bigl(e^{\sigma} - e^{1-\sigma}\bigr)\ln2\;\sin(t\ln2).
\end{align*}
This vanishes iff $\sigma=1/2$ or $t\ln2=n\pi$.
\end{proof}

\section{Conjecture}

Based on the algebraic and numerical evidence, we propose:

\[
\boxed{
\text{If }\zeta(s)=0\text{ and }0<\Re(s)<1,\text{ then }K(s)K(1-s)\in\mathbb{R}.
}
\]

This conjecture is equivalent to the Riemann Hypothesis, since the
phase condition $K(s)K(1-s)\in\mathbb{R}$ forces $\Re(s)=1/2$
(except for the exceptional lines $t=n\pi/\ln2$, which are
numerically excluded but await a rigorous analytic proof).

\section{Numerical Evidence}

We scanned 200 exceptional lines ($n=1,\dots,200$,
$\sigma=0.05,\dots,0.95$, step $0.05$) using mpmath with 30-digit
precision.  No candidate with $|\zeta(s)|<10^{-4}$ was found.
This evidence is consistent with the conjecture but does not
constitute a proof.

\section{Conclusion}

The conjecture is open.  A proof of the remaining step---excluding
the exceptional lines analytically---would prove the Riemann
Hypothesis.

\end{document}
